\subsection{Sonderfälle bei Gleichungen mit Parameter}
\begin{LA}
    {Das Oma-Problem}{Oma}
    Bobs Grossmutter ist \(d\) Jahre älter als seine Mutter. Bob behauptet: ``Wenn ich \(d\) Jahre alt bin, bin ich doppelt so alt wie meine Grossmutter!''

    Das mag seltsam klingen, aber Bob kann das sogar beweisen:

    Sei \(M\) das Alter von Bobs Mutter, \(G\) jenes von Bobs Grossmutter und \(d\) jenes von Bob.

    \textbf{Voraussetzung:}
    \[
        G-M=d
    \]
    oder umgestellt:
    \[
        G=M+d
        \qquad\text{beziehungsweise}\qquad
        M=G-d.
    \]

    \textbf{Behauptung:}
    \[
        d=2G.
    \]

    \textbf{Beweis:}
    Die Reihenfolge bei der Multiplikation spielt bekanntlich keine Rolle, also ist
    \[
        G\cdot M=M\cdot G.
    \]

    Nun verwenden wir auf der linken Seite die Voraussetzung \(M=G-d\) und rechts \(G=M+d\):
    \[
        \begin{aligned}
            G(G-d)    & = M(M+d)     &  & \mid -MG \\
            G(G-d)-MG & = M(M+d)-MG.
        \end{aligned}
    \]

    Auf der linken Seite können wir \(G\) ausklammern, rechts hingegen \(M\). Man erhält
    \[
        \begin{aligned}
            G(G-d-M) & = M(M+d-G)
                     &                           & \mid :(G-d-M) \\
            G        & = \frac{M(M+d-G)}{G-d-M}.
        \end{aligned}
    \]

    Die Klammern im Zähler und im Nenner unterscheiden sich nur in den Vorzeichen:
    \[
        \begin{aligned}
            G
             & =\frac{M(M+d-G)}{G-d-M}           \\
             & =\frac{M\cdot(-1)(-M-d+G)}{G-d-M} \\
             & =\frac{-M(G-d-M)}{G-d-M}          \\
             & =-M.
        \end{aligned}
    \]

    Offenbar gilt also
    \[
        G=-M.
    \]
    Wenn man auf beiden Seiten \(G\) addiert, erhält man
    \[
        2G=G-M=d.
    \]

    Also ist Bob tatsächlich doppelt so alt wie seine Grossmutter!

    \bigskip

    \textbf{Aufgabe: Finden Sie den Fehler im Beweis.}
\end{LA}







\clearpage
\begin{satz}{Das OMA-Problem}{}
    Durch Null darf nicht geteilt werden, sonst lässt sich alles beweisen\footnote{Auch dass man doppelt so alt wie die eigene Grossmutter ist!}.
\end{satz}

Dieser Merksatz ist nicht neu. Durch die Verwendung von Parametern merkt man aber nicht immer gleich, wenn man durch Null teilt!

Eine lineare Gleichung in \(x\) hat nach dem Umformen immer die Form
\[
    a\cdot x=b.
\]

Bevor man mit dem Lösen weiterfährt, muss man an dieser Stelle drei Fälle unterscheiden:

\begin{enumerate}
    \item Falls \(a\neq 0\) ist, kann man durch \(a\) teilen. Die Gleichung kann wie folgt gelöst werden:
          \[
              \begin{aligned}
                  a\cdot x & = b           &  & \mid :a \\
                  x        & = \frac{b}{a}
              \end{aligned}
          \]

    \item Falls \(a=0\) und \(b=0\) ist, hat die Gleichung die Form
          \[
              0\cdot x=0.
          \]
          Diese Gleichung ist allgemeingültig. Jede Zahl ist eine Lösung. Die Lösungsmenge ist
          \[
              \mathbb{L}=\mathbb{R}.
          \]

    \item Falls \(a=0\) und \(b\neq 0\) ist, hat die Gleichung die Form
          \[
              0\cdot x=b.
          \]
          Diese Gleichung ist unerfüllbar. Es gibt keine Lösung. Die Lösungsmenge ist
          \[
              \mathbb{L}=\varnothing.
          \]
\end{enumerate}

Wenn man beim Lösen einer Gleichung mit Parametern die Sonderfälle berücksichtigen will, muss man diese Fallunterscheidung durchführen.
\clearpage
\begin{bsp}{Gleichung mit Sonderfällen}{}
    Die folgende Gleichung soll nach \(x\) aufgelöst werden -- mit Diskussion der Sonderfälle:
    \[
        px=q-5-3x.
    \]
    \vspace*{5cm}
\end{bsp}

\begin{ueb}
    {Diskussion der Sonderfälle}{}
    \begin{enumerate}
        \item Diese Aufgaben haben Sie bereits gelöst, ohne mögliche Sonderfälle zu berücksichten. Lösen Sie sie erneut bzw. ergänzen Sie Ihre Lösung durch die Diskussion der Sonderfälle.

              \begin{enumerate}
                  \item \(cx+c=dx\)
                  \item \(ax=a^2+a\)
                  \item \(qx=1-px\)
                  \item \(qx=-px\)
                  \item \(ax-b=2b+x\)
                  \item \((a+b)x=a^2+2ab+b^2\)
                  \item \(x-p=p^2-px\)
                  \item \(q+3qx=x\)
              \end{enumerate}
    \end{enumerate}
\end{ueb}

\clearpage
\emph{Lösungen:}
\begin{enumerate}
    \item
          \[
              \begin{cases}
                  \mathbb{L}=\left\{\dfrac{c}{d-c}\right\}, & d\neq c,  \\[1ex]
                  \mathbb{L}=\mathbb{R},                    & d=c=0,    \\
                  \mathbb{L}=\varnothing,                   & d=c\neq0.
              \end{cases}
          \]

    \item
          \[
              \begin{cases}
                  \mathbb{L}=\{a+1\},    & a\neq0, \\
                  \mathbb{L}=\mathbb{R}, & a=0.
              \end{cases}
          \]

    \item
          \[
              \begin{cases}
                  \mathbb{L}=\left\{\dfrac{1}{p+q}\right\}, & p+q\neq0, \\[1ex]
                  \mathbb{L}=\varnothing,                   & p+q=0.
              \end{cases}
          \]

    \item
          \[
              \begin{cases}
                  \mathbb{L}=\{0\},      & p+q\neq0, \\
                  \mathbb{L}=\mathbb{R}, & p+q=0.
              \end{cases}
          \]

    \item
          \[
              \begin{cases}
                  \mathbb{L}=\left\{\dfrac{3b}{a-1}\right\}, & a\neq1,       \\[1ex]
                  \mathbb{L}=\mathbb{R},                     & a=1,\ b=0,    \\
                  \mathbb{L}=\varnothing,                    & a=1,\ b\neq0.
              \end{cases}
          \]

    \item
          \[
              \begin{cases}
                  \mathbb{L}=\{a+b\},    & a+b\neq0, \\
                  \mathbb{L}=\mathbb{R}, & a+b=0.
              \end{cases}
          \]

    \item
          \[
              \begin{cases}
                  \mathbb{L}=\{p\},      & p\neq-1, \\
                  \mathbb{L}=\mathbb{R}, & p=-1.
              \end{cases}
          \]

    \item
          \[
              \begin{cases}
                  \mathbb{L}=\left\{\dfrac{q}{1-3q}\right\}, & q\neq\dfrac13, \\[1ex]
                  \mathbb{L}=\varnothing,                    & q=\dfrac13.
              \end{cases}
          \]
\end{enumerate}